\chapter{Tecnologie Utilizzate}
\label{ch:tecnologie}

In questo capitolo vengono descritte le principali tecnologie impiegate nella realizzazione di ChillChain, evidenziando come ciascuna abbia contribuito alla soluzione complessiva.

\section{Spring Boot}

Spring Boot è il framework Java scelto per lo sviluppo del backend della piattaforma, per la sua capacità di semplificare la configurazione e la gestione di applicazioni complesse.

Nel contesto di ChillChain, Spring Boot gestisce le API REST che coordinano il ciclo di vita delle entità di dominio: veicoli, viaggi, spedizioni, notifiche e utenti. Spring Security implementa l'autenticazione tramite token JWT. L'integrazione con MongoDB avviene tramite Spring Data MongoDB, che consente l'utilizzo di repository dichiarativi e pipeline di aggregazione personalizzate. 
Un elemento chiave è l’integrazione con il protocollo MQTT tramite Spring Integration: il backend sottoscrive i topic di telemetria e anomalie pubblicati dai servizi Python, riceve i messaggi tramite channel adapter e li instrada ai service activator responsabili della persistenza dei dati e della generazione delle notifiche.

\section{React e TypeScript}

Il frontend di ChillChain è una SPA (Single Page Application), realizzato con React, una libreria JavaScript basata su componenti, combinata con TypeScript, un superset tipizzato di JavaScript che aggiunge controllo statico dei tipi.

React ha permesso di strutturare l'interfaccia come un insieme di componenti riutilizzabili, ciascuno responsabile di una porzione specifica della vista. La gestione dello stato applicativo si affida agli hook nativi (\texttt{useState}, \texttt{useEffect}) e a \texttt{localStorage} per la persistenza delle informazioni di sessione. La visualizzazione dei dati telemetrici in tempo reale utilizza la libreria Recharts, mentre le tratte dei veicoli vengono mostrate su mappa interattiva tramite React-Leaflet. Il progetto frontend è costruito con Vite e stilizzato con Tailwind CSS, garantendo un’interfaccia coerente e responsive.

\section{MQTT e Mosquitto}

MQTT (\textit{Message Queuing Telemetry Transport}) è il protocollo leggero basato sul pattern \textit{publish/subscribe} adottato per lo scambio dei dati telemetrici tra i servizi della piattaforma.

In ChillChain, MQTT connette tre attori principali: il servizio di streaming dei dati (publisher), il servizio di anomaly detection (subscriber e publisher) e il backend Spring Boot (subscriber). Il broker Mosquitto gestisce l'instradamento dei messaggi tramite topic.
\section{MongoDB}

MongoDB è il database NoSQL orientato ai documenti impiegato per la persistenza di tutti i dati della piattaforma.

Nel contesto di ChillChain, la flessibilità di MongoDB si è rivelata particolarmente vantaggiosa per la gestione dei dati telemetrici, la cui struttura può variare in funzione del tipo di sensore, e per le entità di dominio come viaggi e spedizioni, che includono campi con struttura annidata (coordinate geografiche, polyline).  L'interazione avviene tramite Spring Data MongoDB, che mappa automaticamente le classi Java alle collezioni corrispondenti. Inoltre, vengono utilizzate le pipeline di aggregazione di MongoDB per implementare la logica di ricerca dei viaggi già programmati e per la produzione di report.

\section{Python, FastAPI e PyTorch}

Python è il linguaggio utilizzato per i due microservizi dedicati rispettivamente allo streaming dei dati telemetrici e alla rilevazione delle anomalie. La scelta è motivata dalla ricchezza dell'ecosistema Python per il \textit{data processing} e il \textit{machine learning}.

Il servizio di streaming (\texttt{streamer\_main.py}) è un'applicazione FastAPI che legge i file CSV contenenti i dati simulati di telemetria e li pubblica sul broker MQTT con una cadenza configurabile, simulando il comportamento di sensori reali installati a bordo dei veicoli. I dataset sono generati da uno script di simulazione termodinamica, derivato dal dataset Kaggle \textit{Simulated Refrigerator Fault Diagnosis} \cite{kaggle_fridge}, che modella il comportamento fisico di un impianto frigorifero.

Il servizio di anomaly detection (\texttt{anomaly\_main.py}) è anch'esso un'applicazione FastAPI che sottoscrive il topic telemetria via MQTT, analizza i dati in ingresso e pubblica i risultati sul topic anomalie. Il cuore del servizio è un modello LSTM Autoencoder implementato in PyTorch: l'encoder, una rete LSTM, comprime la sequenza temporale di ingresso in una rappresentazione latente a dimensionalità ridotta; il decoder, un'altra rete LSTM, ricostruisce la sequenza originale a partire da tale rappresentazione. L'errore di ricostruzione, calcolato come errore quadratico medio (MSE) tra input e output, viene confrontato con una soglia appresa durante l'addestramento su dati di funzionamento normale. Lo scaler, serializzato con Pickle, normalizza le quindici feature in ingresso per garantire che parametri con scale numeriche molto diverse contribuiscano in modo equilibrato all'errore di ricostruzione. La libreria Pandas è impiegata per la lettura dei CSV e la preparazione dei DataFrame, mentre NumPy gestisce le operazioni numeriche sui buffer a finestra scorrevole.

\section{Google Maps Platform}

Le API di Google Maps Platform sono utilizzate per il calcolo dei percorsi, la geolocalizzazione degli indirizzi e l'assistenza all'inserimento degli stessi da parte degli utenti. Il sistema impiega tre servizi distinti, distribuiti tra frontend e backend.

Sul lato frontend, il componente di inserimento spedizione integra la \textit{Places Autocomplete API}, che fornisce suggerimenti in tempo reale durante la digitazione degli indirizzi di ritiro e consegna. Oltre a migliorare l'esperienza utente riducendo errori di battitura e ambiguità, il servizio restituisce direttamente le coordinate geografiche associate all'indirizzo selezionato, eliminando la necessità di una successiva chiamata di geocodifica.

Sul lato backend vengono utilizzate due API aggiuntive: la \textit{Geocoding API}, che converte gli indirizzi testuali in coordinate geografiche (latitudine e longitudine) nei casi in cui le coordinate non siano già disponibili dal frontend, e la \textit{Routes API}, che calcola il percorso ottimale tra due o più punti restituendo la distanza in metri, la durata stimata e la polyline codificata del tracciato.

\section{Docker e Docker Compose}

Docker è la piattaforma di containerizzazione impiegata per isolare ciascun componente dell'architettura in un ambiente di esecuzione autonomo e riproducibile. Ogni servizio della piattaforma viene eseguito all'interno di un container dedicato, con le proprie dipendenze e la propria configurazione.

Docker Compose orchestra l'avvio coordinato di tutti i container, definendo in un unico file YAML le immagini, le porte esposte, le variabili d'ambiente, i volumi persistenti e le reti interne. 